% !TEX TS-program = lualatex

%%%
% src::
%     url = https://tex.stackexchange.com/a/750411/6880
%%%

\documentclass[border = 3pt]{standalone}

%%%
% Des couleurs faciles d'emploi via le package \xcolor.
%%%
\usepackage[svgnames]{xcolor}

%%%
% La bibliothèque \luadraw allie une facilité d’utilisation
% à un rendu particulièrement soigné.
%%%
\usepackage[3d]{luadraw}

\begin{document}

\begin{luadraw*}{name = implicit-tangent-lines}
local ld  = luadraw
local Z   = ld.cpx.Z
local abs = ld.cpx.abs

------
-- L'équation implicite est `F(x, y) = 0`.
------
local F = function(x, y)
  return x^2 + y^4 + x * y - 1
end

------
-- ¨Def de la zone graphique.
------
local graphview = ld.graph:new{
  window = {-3, 3, -3, 3},
  size   = {10, 10},
  bbox   = false
}

------
-- Il est temps de dessiner...
------
graphview:Dgradbox(
  {
    Z(-2, -2), Z(2, 2),
    1, 1
  },
  {
    grid  = true,
    title = "{\\boldmath $x^2 + y^4 + x y = 1$}"
  }
)

------
-- La valeur par défaut de ''grid'' rend visible les segments
-- constituant la courbe, donc nous affinons les pas de calcul.
------
graphview:Dimplicit(
  F,
  {
    grid         = {100, 100},
    draw_options = "thick"
  }
)

------
-- Ajout de tangentes aux points d'abscisses `0.5`.
------
local x0 = 0.5

local L = ld.solve(
  function(t)
    return F(x0, t)
  end,
  -2, 2
)

for _,  y in ipairs(L) do
    graphview:DtangentI(
        F, x0, y, 2,
        "Stealth-Stealth, thick, red"
    )
end

------
-- Montrons le résultat de notre oeuvre.
------
graphview:Show()
\end{luadraw*}

\end{document}
